\subsection{Relation to RACs and  MI-RNNs}
\label{appendix: sec: claim 4.3}

We here focus on Multiplicative Integration (MI), a way to connect two sources of inputs by the Hadamard product `$\odot$'. MI has been used in RACs, Multiplicative RNNs (M-RNNs) \citep{sutskever2011generating} and MI-RNNs: 

\emph{Recurrent Arithmetic Circuits} (RACs) are recurrent networks with  hidden states $\vh^{(t)}_\textrm{RAC}$ defined as  \citep{levine2018benefits}:
%
\begin{equation} \small
\label{eq: RAC_simple}
   \vh^{(t)}_\textrm{RAC} =  \mW^{hx} \vf (x^{(t)}) \odot \mW^{hh}\vh^{(t-1)}_\textrm{RAC}
\end{equation}
%

where these hidden states are also used as an \emph{intermediate term} in M-RNNs. 

\emph{Multiplicative Integration RNNs} (MI-RNNs) are RACs with an activation function and hidden states $\vh^{(t)}_\textrm{MI}$ defined as  \citep{wu_multiplicative_2016}:


\begin{equation} \small
\label{eq: mirnn}
    \vh^{(t)}_\textrm{MI}= \phi(\mW^{hx} \vf(x^{(t)}) \odot \mW^{hh}\vh^{(t-1)}_\textrm{MI})
\end{equation}

% \begin{wrapfigure}[16]{r}{0.42\textwidth}
% \small
% % \vspace{-10mm}
% \begin{center}
% \includegraphics[width=0.40\textwidth]{figure/SecondOrderRNN.pdf}
% \vspace{-4mm}
% \caption{\small{Second-order RNNs}}.
% \vspace{-2mm}
% \label{fig:secondorderrnn}
% \end{center}
% \end{wrapfigure} 




% \subsection{TTLM generalizes Second-order RNNs }
% \label{sec: relationship}

\begin{figure}[t]
    \centering
    \includegraphics[scale=0.45]{figure/TTLM_Tiny.pdf}
    \caption{a) Second-order RNNs under TTLM framework.  b) Hidden state of RACs under TTLM framework. c) hidden state of MI-RNNs under TTLM framework. The dashed line in the square denotes $\mathcal{A}, \Phi(X)$ or $\tG$. The small hollow circles denote the activation functions.}
  \label{fig:rnns}   
\end{figure}



% \begin{proof} The detailed proof can be found in Appendix \ref{appendix: seecond}. The diagrammatic notation of Second-order RNNs is shown in Fig. \ref{fig: TTLM}b.

% \end{proof}





% \begin{wrapfigure}[16]{r}{0.42\textwidth}
% \small
% % \vspace{-10mm}
% \begin{center}
% \includegraphics[width=0.40\textwidth]{figure/RACs.pdf}
% \vspace{-4mm}
% \caption{\small{RACs}}.
% \vspace{-2mm}
% \label{fig:RACs}
% \end{center}
% \end{wrapfigure} 




\begin{claim}
\label{cm: G and RACs} Given the condition the TT-scores: $\tG = \mW^{hx} \odot \mW^{hh}$. 
The hidden states of RACs are identical to that of TTLM. There is a nonlinear function $\phi$ such that the hidden states of MI-RNNs are identical to that of TTLM if they are accompanied by $\phi$.   
\end{claim}

% \begin{proof} The detailed proof can be found in Appendix \ref{appendix: sec: claim 4.3}. The diagrammatic notations of RACs and MI-RNN  are shown in Fig. \ref{fig: TTLM}.
% \end{proof}


\begin{proof} The proof is based on the following observation:
%
In the language of tensor contractions, Eq.\ \ref{eq: RAC_simple} involves contracting the input weights matrix  $\mW^{hx}$ with the input vector $\vf (x^{(t)})$, and contracting the hidden weights matrix $\mW^{hh}$ with $\vh^{(t-1)}_\textrm{RAC}$. The Hadamard product of the two is a third-order diagonal tensor $\delta \in \mathbb{R}^{R \times R \times R}$ such that $\delta_{ijk} = 1$ if{f} the $i = j = k$, and $\delta_{ijk} = 0$ otherwise. Thus, we can write Eq.\ \ref{eq: RAC_simple} in element-wise fashion: 
%
\begin{equation} \small
\begin{split}
    \label{eq: RAC_element}
    h^{(t)}_{\textrm{RAC}_{\alpha_t}}
    & = \sum_{{i_t = 1}}^{|V|} \sum_{{\alpha_t = 1}}^{R} f(x^{(t)})_{i_t}W^{hx}_{i_tj}\mathbf{\delta}_{j\alpha_{t}k}W^{hh}_{k\alpha_{t-1}} h^{(t-1)}_{\textrm{RAC}_{\alpha_{t-1}}} \\
    & = \sum_{{i_t = 1}}^{|V|} \sum_{{\alpha_t = 1}}^{R} f(x^{(t)})_{i_t}\etG_{i_t\alpha_{t}\alpha_{t-1}} h^{(t-1)}_{\textrm{RAC}_{\alpha_{t-1}}} \\
\end{split}
\end{equation}


where $\tG = \mW^{hx} \odot \mW^{hh}$. In this case, the hidden state of TTLM in Eq. \ref{eq: ttlm cell} is equal to the hidden state of RACs in Eq. \ref{eq: RAC_element},  as shown in Fig. \ref{fig:rnns}b.  Similarly,  if Eq. \ref{eq: ttlm cell} is accompanied with an activation function $\phi$,  Eq. \ref{eq: ttlm cell}  is equal to the hidden state of MI-RNNs in Eq. \ref{eq: mirnn} as shown in Fig. \ref{fig:rnns}c. 

\end{proof}

Given Claim \ref{cm: G and T} and \ref{cm: G and RACs}, the three models shall be simulated by TTLM with a nonlinear activation function and we leave finding a theoretical proof of this conjecture to a future work.

% \section{Implementations}
% \label{sec:implementations}
% % Note that in order to follow a realistic case, the baseline models implemented by us have an embedding matrix $W^{xe}$ between one-hot vectors $\vf (x^{(t)})$ and the input-to-hidden  matrix $\mW^{hx}$ \footnote{This change  obeys the multiplicative form and follows the TT-format of TTLM ($W^{xe}$ is part of TT cores)}. 
% We implement all RNNs models, TTLM, TTLM-Large, and TTLM-Tiny using PyTorch on GPU A100 one single card.


% For RNNs, there are five matrix parameters: $\mW^{xe}\in \mathbb{R}^{E\times |V|}$ is the input embedding matrix, $\mW^{eh}\in
% \mathbb{R}^{E\times H}$ is the embedding-to- hidden matrix, $\mW^{hh}\in \mathbb{R}^{H\times H}$ is the hidden-to-hidden matrix. We tie (share the same training parameters) the input embedding  $\mW^{xe}$ and output embedding $\mV$ which has been proved lead to a significant reduction in perplexity \citep{press2016using}. So there is a projection matrix $\mP\in\mathbb{R}^{H\times E}$ before the output embedding. All this process is introduced in  \citep{press2016using}. 

% For TTLM models, we tie the input tensor $\tW^{xe}$ and $\tV$. The implementation of $\delta$ is functioned by a reshape function, so the interaction between hidden and input can be computed by matrix product. We also let $\tG^{(1)}$ have the same parameters along the dimension $|V|$ (i.e. $\tG^{(1)}$ is simplified into a  $\mG^{(1)}\in\mathbb{R}^{1\times R}$ and thus it can be viewed as the initial hidden state $\vh^{(0)}$).  



% \begin{table}[ht]
%     \centering
%     \begin{tabular}{l|c}
%     \hline
%         Model & Training Parameters \\
%         \hline
%         Vanilla-RNN & $\mW^{xe}\in \mathbb{R}^{E\times |V|}$, $\mW^{eh}\in \mathbb{R}^{E\times H}$, $\mW^{hh}\in \mathbb{R}^{H\times H}$, \\
%         & $\mP\in\mathbb{R}^{H\times E}$, $\mV\in\mathbb{R}^{E\times|V|}$  \\
%         \hline
        
%         MI-RNNs & $\mW^{xe}\in \mathbb{R}^{E\times |V|}$, $\mW^{eh}\in \mathbb{R}^{E\times H}$, $\mW^{hh}\in \mathbb{R}^{H\times H}$, \\ & $\mP\in\mathbb{R}^{H\times E}$, $\mV\in\mathbb{R}^{E\times|V|}$\\
%         \hline
%         RACs & $\mW^{xe}\in \mathbb{R}^{E\times |V|}$, $\mW^{eh}\in \mathbb{R}^{E\times H}$, $\mW^{hh}\in \mathbb{R}^{H\times H}$, \\
%         & $\mP\in\mathbb{R}^{H\times E}$, $\mV\in\mathbb{R}^{E\times|V|}$\\
%         \hline
%         Second-order RNNs & $\mW^{xe} \in \mathbb{R}^{E\times |V|}$, $\tT \in \mathbb{R}^{H\times H \times H}$, $\mW^{hh}\in \mathbb{R}^{E \times H}$, \\
%         & $\mP\in\mathbb{R}^{H\times E}$, $\mV\in\mathbb{R}^{E\times|V|}$\\
%         \hline
%         TTLM & $\tG\in \mathbb{R}^{R\times |V|\times R}$, $\mG^{(t)}\in \mathbb{R}^{R \times |V|}$, $\mG^{(1)} \in \mathbb{R}^{R \times |V|}$ \\
%         \hline
%         TTLM-Tiny & $\tW^{xe}\in \mathbb{R}^{R\times |V| \times R}$, $\mW^{hh}\in\mathbb{R}^{R\times R}$, \\
%         & $\tP\in\mathbb{R}^{R\times R\times R}$, $\tV\in\mathbb{R}^{R\times R\times |V|}$\\
        
%         \hline
%         TTLM-Large & $\tW^{xe}\in\mathbb{R}^{R\times |V|\times R}$, $\tW^{eh}\in\mathbb{R}^{R\times R\times R\times R}$, $\mW^{hh}\in\mathbb{R}^{R\times R}$, \\
%         & $\tP\in\mathbb{R}^{R\times R\times R}$, $\tV\in\mathbb{R}^{R\times R\times |V|}$\\
%     \end{tabular}
%     \caption{Training parameters in our implementation. $E$ is the embedding size, $H$ is the hidden units in RNNs, and $R$ is the rank in the TTLM.  We set $H=R$ and $E=R^2$ to make the parameters of all models in the same scale. The parameters of $\mW^{xe}$  and $\tW^{xe}$are uniformly initialized in the interval $[-0.1,0.1]$, $\mW^{eh}$, $\tW^{eh}$ and $\mW^{hh}$ are uniformly initialized between $[-\frac{1}{\sqrt{H}}, \frac{1}{\sqrt{H}}]$.}
%     \label{tab:my_label}
% \end{table}
% %%%%%%%%%%%%%%%%%%%%%%%%%%%%%%%%%%%%%%%%%%%%%%%%%%%%%%%%%%%%%%%%%%%%%%%%%%%%%%%
% %%%%%%%%%%%%%%%%%%%%%%%%%%%%%%%%%%%%%%%%%%%%%%%%%%%%%%%%%%%%%%%%%%%%%%%%%%%%%%%
